\chapter{结论}
\label{ch:conclusion}

\section{总结}

本文研究如何在 Godot 4.6 中构建完整可视化行为树，同时使大型树更容易检查。项目面对两个相互连接的问题：Godot 提供通用图编辑接口，却没有完整行为树流程；传统完整细节节点图会随着节点、参数和运行标注增加而拥挤。

最终 0.9.0 插件实现可序列化树、类型化黑板、Decorator、有序与响应式复合节点、随机与并行控制、Repeat/Wait 记忆、Actor 方法 Action/Condition、可复用 NPC 组件、撤销重做和仅编辑器 Live Debug。复杂二维场景使用 241 节点资源展示索敌、响应式防御、近战与远程攻击、越障跳跃、梯子追击、最后位置搜索、撤退、治疗和巡逻。

RQ1 在定义范围内得到肯定回答：Godot 编辑器与 Resource API 可以支撑实用的编写、执行和调试流程。RQ2 的仪器化部分同时由控制实验和可玩树证据回答：在 241 节点资源上，Compact 把卡片面积降低 57.93\%，Overview 降低 89.48\%，Search 淡化 202 张卡片中的 201 张，Focus/Collapse 将画布缩减至 32/43 张卡片；所有条件保持零重叠和保存位置，固定路径拖拽处理降低 26.60\%。RQ2 的人体任务表现部分尚未回答：平衡的 12 人/216 次实验已经实现，但当前观测为零。RQ3 在当前工作负载下得到肯定回答：精确签名缓存保持执行结果，并将 Tick 中位成本降低 15.98--78.00\%，大树收益更高。

\section{贡献}

本项目贡献了可用 Godot 插件，而不是脱离运行时的视觉原型；贡献了集成的可切换大树与运行检查机制、安全拓扑缓存，以及包含自动回归、180 条可玩树原始显示观测、同任务拖拽比较、固定截图和已校验开发者实验工作簿的可复现证据包。它也提供一个重要的方法论结论：几何压缩不能直接报告为人的可读性已经改善；预注册协议和空数据保护使这条边界保持明确。

\section{局限}

当前证据限于一台 Windows/GPU、Godot 4.6、最多 364 节点的生成树和一个可玩 241 节点工作负载。自动结果能证明实现与表示属性，却不能证明长期编写效率或认知负荷。调试桥提供最新状态而非完整时间线。Service、可复用子树资源、Asset Library 和跨版本矩阵不在已完成范围内。

\section{未来工作}

最优先工作是按已准备协议招募 12 人，在 216 次平衡试验中比较 Baseline、Compact、Collapse、Fisheye、Search 和 Minimap 的 Action 定位、活动路径和 Decorator 编辑任务。上限完成时间、成功率、错误、平移缩放和主观评分按任务分别采用重复测量检验、校正比较和效应量分析，且不能把不可直接比较的几何条件合并成无依据的总分。

工程上可增加可复用子树接口、暂停/单步调试时间线、多 Runner Actor 选择，以及“事件版本号+周期精确校验”的混合缓存。显示研究可以加入增量动画与可配置兴趣度函数，并测量重访和心理地图。最后还应验证选定 Godot 版本并准备 Asset Library 上架。

\section{最终结论}

本项目说明毕业设计可以超越对既有编辑器的简单模仿。通过结合可执行语义、仅编辑器解释、可切换焦点与复杂度控制、严格回归和明确证据边界，该插件既是实用 AI 编写工具，也是后续研究可读性可视化编程的平台。
